\makeabstract
    {
    Optical Accretion Signatures From Planetary Mass Companion: Delorme 1 (AB)b
    }
    {
    Valentina Guerrero Cachucho\textsuperscript{1}, Katherine Follette\textsuperscript{1}, Rebecca Michelson\textsuperscript{1}
    }
    {
    \textsuperscript{1} Department of Physics \& Astronomy, Amherst College, Amherst, MA
    }
    {
    Accretion is the process by which an object accumulates or accretes mass, usually driven by gravitational or magnetic forces. This process is especially intertwined with the earlier stages of formation in stars, brown dwarfs, and planets, as it represents the period in which these objects are actively growing. It is by observing accretion that we can better understand the mechanisms by which these objects form. While the physical mechanisms of accretion in stars are well-studied, this paradigm may not apply to lower-mass objects like brown dwarfs and planets. The aim of this thesis is to explore the signatures of accretion, like hydrogen line emission and continuum excesses, in the optical spectrum of the planetary mass companion (PMC), Delorme 1 AB b. By investigating an accreting PMC, we can gain insight into the processes of planetary accretion and thus the formation and growth of gaseous planets.
    }
    {
    Optical spectra were obtained from Cycle 30 of the Hubble Space Telescope (HST) using the Space Telescope Imaging Spectrograph (STIS). These optical spectra were manually re-extracted from 2D spectroscopic images as the automatic STIS pipeline extracted the spectrum at the wrong location. The extracted spectra were then combined using a few different methods (e.g. inverse-variance weighting, adding fluxes) to increase the spectrum signal-to-noise ratio. Spectra from JWST/MIRI, TripleSpec, and VLT/MUSE in the infrared wavelengths were also acquired courtesy of Betti et al. 2022 and Malin et al. 2025. These observations were plotted alongside the HST data, placing the novel HST/STIS data into context with past available data.
    }
    {
    The emission lines from the re-extracted Delorme 1 AB b spectrum masked for data quality flags look as expected, with the exception of the H-gamma line, which has been flagged as a hot pixel at its peak. Further data processing will be conducted to adjust the bad pixel threshold to avoid excluding the H-gamma emission. Moreover, the continuum in the optical range seems to be undetected, likely due to emitted fluxes below the STIS detection limit. 
    }
    {
    Next steps include fitting a planetary atmospheric model to the infrared observations in order to extrapolate the model back into the optical wavelengths and get a continuum level estimate. The emission line fluxes in the HST spectrum will then be compared to the model continuum, allowing for the measurement of line fluxes and line luminosities. Then, using the empirical stellar and theoretical planetary Lacc-Lline relations, the mass accretion rate for Delorme 1 Ab b can be estimated. Ultimately, our goal is to measure mass accretion rates using H emission lines and compare the mass accretion rates measured using the empirical stellar Lacc-Lline relations from Alcala et al. 2017 and the theoretical Lacc-Lline relations from Aoyama et al. 2021. This will add 6 new individual H-line mass accretion rate measurements to the literature values from Betti et al. 2022.
    }

    \newpage
    